% pdflatex file_name.tex
\documentclass{resume} % Use the custom resume.cls style
\usepackage{textcomp}
\usepackage[left=0.5in,top=0.5in,right=0.5in,bottom=0.5in]{geometry} % Document margins
\newcommand{\tab}[1]{\hspace{.2667\textwidth}\rlap{#1}}
\newcommand{\itab}[1]{\hspace{0em}\rlap{#1}}
\usepackage{ragged2e}
\justifying
\usepackage{hyperref}
\usepackage{multicol}
\usepackage{graphicx}
\usepackage{pifont}
\usepackage{enumitem}
\usepackage{array}
\usepackage{multicol}
\usepackage{acronym}
\newacro{DECam}[DECam]{Dark Energy Camera}
\newacro{desc}[DESC]{Dark Energy Science Collaboration}
\newacro{LSSTCam}[LSSTCam]{LSST-Camera}
\newacro{ToO}[ToO]{Target-of-Opportunity}
\newacro{igwn}[IGWN]{International Gravitational-Wave Observatory Network}
\newacro{desgw}[DESGW]{Dark Energy Survey gravitational wave}
\newacro{O4}[O4]{observing run 4}
\newacro{J-GEM}[J-GEM]{Japanese collaboration for Gravitational wave ElectroMagnetic follow-up}
\newacro{GW}[GW]{gravitational wave}
\newacro{PFS}[PFS]{Prime Focus Spectrograph}
\newacro{DELVE}[DELVE]{DECam Local Volume Exploration survey}
\newacro{osi}[PSI]{Paul Scherrer Institute}
\newacro{des}[DES]{Dark Energy Survey}
\newacro{LLAMAS}[LLAMAS]{Large Lenslet Array Magellan Spectrograph}
\newacro{NERSC}[NERSC]{National Energy Research Scientific Computing Center}
\newacro{lsst}[LSST]{Legacy Survey of Space and Time}


\renewcommand\labelitemi{\ding{226}}
\name{Sean MacBride} % Your name
% \address{40 Wildwood Drive \\ Burlington, VT 05408} % Your address
%\address{123 Pleasant Lane \\ City, State 12345} % Your secondary addess (optional)
\address{\href{mailto:Sean.P.MacBride@gmail.com}{Email} \\ \href{https://github.com/seanmacb}{GitHub} \\ \href{https://seanmacb.github.io}{\texttt{seanmacb.github.io}} \\ \href{tel:+18029224673}{+1 (802) 922 4673}} %\\ \href{https://www.linkedin.com/in/sean-macbride/}{LinkedIn}}% Your phone number and email
% Add website when it's finished
\usepackage{seqsplit}
\begin{document}
\begin{rSection}{Research Interests}
% \begin{itemize}
%     \item Experimental and Observational Cosmology
%     \item Astronomical Instrumentation Development
%     \item Multi-Messenger Astrophysics
%     \item Systems Design for Astronomy
% \end{itemize}

\begin{center}    
\begin{tabular}{ll}
$\bullet$ Experimental and Observational Cosmology & $\bullet$ Astronomical Instrumentation Development \\
$\bullet$ Multi-Messenger Astrophysics             & $\bullet$ Systems Design for Physics \& Astronomy            
\end{tabular}
\end{center}
\end{rSection}



% \begin{rSection}{Honors and Awards}
% \begin{tabular}{ll}
%     2018 & Boggess Family Foundation Scholarship \\
%     2017 & Wheaton Centennial Grant
% \end{tabular}    
% \end{rSection}

\begin{rSection}{Scientific Collaborations}

\begin{rSubsection}{Rubin Observatory \ac{lsst}}{2023 - Present}{}{}
I hold roles in both the \href{https://rubinobservatory.org}{Rubin Project} and the \href{https://lsstdesc.slac.stanford.edu}{\ac{lsst} \ac{desc}}.

\textit{Rubin Observatory Project Team}: Led \ac{LSSTCam} focal plane commissioning to Rubin Observatory in Cerro-Pachón in Chile, from receipt on summit through early science operations, including testing of the focal plane in-laboratory and on-telescope. Characterized defects in the \ac{LSSTCam} focal plane CCDs, and developed custom image processing tool to monitor the stability of defects in the \ac{LSSTCam} focal plane. Led the \ac{ToO} observing program during Rubin early operations and first year of the \ac{lsst}.

\textit{\ac{lsst} \ac{desc}}: Forecasted the constraining power of dark sirens using simulated \ac{lsst} data under different conditions. Leading a dark-siren analysis using \ac{lsst} commissioning data (Data-Preview 2) and public data from the International Gravitational Wave Network (IGWN), the first dark siren measurement using \ac{lsst} data.
\end{rSubsection}

\begin{rSubsection}{\ac{igwn}}{2025 - Present}{}{}

Led dark siren analysis using \ac{DELVE} survey data for \ac{igwn} O4c data, including galaxy catalog preparation and systematic-uncertainty quantification. Coordinated time-critical \ac{ToO} observations with \ac{igwn} researchers, interfacing with Rubin Observatory alert brokers to rapidly disseminate information about \ac{GW} candidates. Contributed to dark siren analysis using Dark Energy Survey (DES) data for \ac{igwn} O4b data.

\end{rSubsection}

\begin{rSubsection}{\href{https://www.darkenergysurvey.org/}{\ac{des}}}{2023 - Present}{}{}
Performed observations, image processing, and analysis of images from \ac{DECam} observations in response to \ac{GW} triggers from the \ac{igwn} \ac{O4} for the \ac{desgw} observing program. Implemented improvements to \ac{desgw} alert and image processing pipelines, including improvements to the telescope strategy and \ac{GW} alert processing. Led collaboration between the \ac{J-GEM} and \ac{desgw} through joint observations of \ac{GW} event S250328ae using \ac{DECam} and the \ac{PFS} at Subaru Observatory.

\end{rSubsection}


\begin{rSubsection}{Robotic Positioners for Future Spectroscopic Surveys}{2023 - Present}{}{}
Tested Dark Energy Spectroscopic Instrument (DESI) robotic positioners to increase movement precision and improve low performing robotic positioners on the DESI focal plane. Characterized a common failure mode of DESI positioners and resolved the failure with a reproducible solution. Authored and submitted an internal report and presented the results of this study to the DESI focal plane working group. Led dynamic testing of prototype robotic positioner performance through lifetime, focal-ratio degradation, and fiber angle tests using different telescope configurations. Tested prototype robotic positioners for Stage-V surveys, focusing on characterization of tilt induced by different orientations.    
\end{rSubsection}

\begin{rSubsection}{\href{https://ait.mit.edu/instruments/llamas/}{\ac{LLAMAS}}}{2020 - 2022}{}{}
Assembled the opto-mechanical mounts for the \ac{LLAMAS} instrument spectrograph. Led verification and testing of mechanical tolerance for custom-fabricated parts with a coordinate measurement machine. Verified spectrograph grating efficiency through optical tests, and integrated the spectrograph camera lenses into their bezels for the 24 camera spectrograph. Integrated the 2,400 AR-coated fibers to the spectrograph by bonding with optical adhesive, with 100\% accuracy and zero fiber loss.
\end{rSubsection}


\end{rSection}

\begin{rSection}{Education}
\begin{rSubsection}
        {Ph.D. in Physics, University of Zürich}{May 2024 - 2027}{}{}
        Thesis title: \textit{Standard Siren Cosmology: From Instrumentation to Inference}
        \newline
        Advisor: Professor Marcelle Soares-Santos
\end{rSubsection}
    
\begin{rSubsection}
        {M.Sc. in Physics, University of Michigan-Ann Arbor}{Aug 2022 - May 2024}{}{}
        Advisor: Professor Marcelle Soares-Santos
\end{rSubsection}

\begin{rSubsection}{B.A. in Physics with Honors, Wheaton College, Massachusetts}{Aug 2016 - May 2020}{}{}

    Thesis title: \textit{\href{https://wheatoncollege.as.atlas-sys.com/repositories/2/archival_objects/21957}{Characterizing the Dust and Cold-Gas Content of Nearby Star-Forming Galaxies.}}
    \newline
    Advisor: Professor Dipankar Maitra

\end{rSubsection}

\end{rSection}

\begin{rSection}{Professional Appointments, Service, and Leadership}
\begin{table}[h]
    \setlength{\extrarowheight}{3pt} 
\begin{tabular}{p{0.1\linewidth}p{0.8\linewidth}}
     2026 - 2028 & \ac{lsst} \ac{desc}, Speakers Bureau Co-Chair \\
     2026 - 2028 & \ac{lsst} \ac{desc}, Collaboration Council Member \\
     2025 - 2027 & Rubin Observatory, \ac{ToO} Observer \\
     2024 - 2027 & Rubin Observatory, \ac{LSSTCam} Summit Support Scientist \\
     2022 - 2024 & University of Michigan Physics Graduate Council, Elected member \\
     2022 - 2024 & University of Michigan Physics DEI committee, Member\\
     2022 - 2024 & University of Michigan APS chapter, Member\\
     2016 - 2020 & Wheaton College physics club, Member
\end{tabular}
\end{table}
\end{rSection}

\begin{rSection}{Teaching}
\begin{table}[h]
    \setlength{\extrarowheight}{3pt} 
\begin{tabular}{p{0.1\linewidth}p{0.8\linewidth}}
2026 & Introduction to Scientific Computing (PHY124), University of Zürich, Graduate Instructor\\
2025 & Introduction to Astroparticle Physics (PHY473), University of Zürich, Graduate Instructor\\
2023 - 2024 & Introductory Physics for Life Sciences (PHYSICS 151/251), University of Michigan-Ann Arbor, Lead Graduate Instructor\\
2022 - 2023 & Introductory Physics for Life Sciences (PHYSICS 151/251), University of Michigan-Ann Arbor, Graduate Instructor\\
2018 - 2019 & Introductory Physics (PHYS 170/171/180/181), Wheaton College, MA, Undergraduate Teaching Assistant and Tutor\\
\end{tabular}
\end{table}    
\end{rSection}

\begin{rSection}{Approved Telescope Programs}
    \begin{table}[h]
    \setlength{\extrarowheight}{3pt} 
\begin{tabular}{p{0.1\linewidth}p{0.8\linewidth}}
        2025 & GW-Primo. PI: McMahon, Co-I: MacBride. 2 nights on the Blanco telescope at CTIO. \\
        2023 - 2027 & DESGW. PI: Soares-Santos, Co-I: MacBride. 4 nights on the Blanco telescope at CTIO, every semester (6 months).
    \end{tabular}
    \end{table}
\end{rSection}
\newpage
\begin{rSection}{Research Grants}
\begin{table}[h]
    \setlength{\extrarowheight}{3pt} 
\begin{tabular}{p{0.1\linewidth}p{0.8\linewidth}}
    2024 - 2026 & \textit{Quantifying the Impact of Galaxy Catalog Conditions for Dark Sirens}, 3M CPU hour grant National Energy Research Scientific Computing Center, ScienceAtScale
\end{tabular}    
\end{table}
\end{rSection}

\begin{rSection}{Conference Organization}
\begin{table}[h]
    \setlength{\extrarowheight}{3pt} 
\begin{tabular}{p{0.1\linewidth}p{0.8\linewidth}}
     2025 &  EAS Annual Meeting Special Session 27: The Rubin Observatory Legacy Survey of Space and Time: from Data Previews to Full Survey Operations, co-chair, local organizing committee \\
     2024 &  Conferences for Undergraduate Women in Physics (CUWiP) 2024, local organizing committee \\
     2023 &  DESI Collaboration Meeting, local organizing committee \\
     2023 & University of Michigan-Ann Arbor Physics Graduate Student Symposium, chair local organizing committee
\end{tabular}
\end{table}
\end{rSection}

\begin{rSection}{Mentorship}
\begin{table}[h]
    \setlength{\extrarowheight}{3pt} 
\begin{tabular}{p{0.1\linewidth}p{0.8\linewidth}}
    2026 & Alexs Gaute: Summer Masters Student at University of Zürich.\\
    2025 - 2026 & Guandi Zhao: Masters Student at University of Zürich. \\
    2024 - 2025 & Banan Yamani: Masters Student at University of Zürich. Current Physics PhD student at Paul Scherrer Institute.\\
    2024 - 2025 & Yavuz Gencel: Masters Student at University of Zürich. \\
    2023 - 2024 & Elise Fuller: Undergraduate student at University of Michigan-Ann Arbor. Current physics PhD student at CU Boulder.\\
    2023 - 2024 & Tim Fanning: Undergraduate student at University of Michigan-Ann Arbor. Current nuclear physics PhD student at Indiana University Bloomington.\\
\end{tabular}
\end{table}
\end{rSection}

% \begin{rSection}{Diversity Equity and Inclusion}
    
% \end{rSection}

\begin{rSection}{Selected Publications}

Below are selected publications from different collaborations and projects. A full list of my publications can be found on my \href{https://scholar.google.com/citations?user=XtRTswUAAAAJ&hl=en}{Google Scholar profile} or on \href{https://ui.adsabs.harvard.edu/search/fq=\%7B!type\%3Daqp\%20v\%3D\%24fq_database\%7D&fq_database=(database\%3Aastronomy\%20OR\%20database\%3Aphysics)&p_=0&q=((\%20author\%3A\%22Macbride\%2C\%20Sean\%22)\%20AND\%20year\%3A2019-2150)&sort=date\%20desc\%2C\%20bibcode\%20desc}{NASA ADS}. My h-index as of \today, is 7.

\begin{rSubsection}{Astronomical Instrumentation}{}{}{}
\begin{enumerate}
    \item \href{https://arxiv.org/abs/2606.31806}{The On-Sky Performance of the LSST Camera CCD Array. \textbf{MacBride}, Roodman, ... 30 Jun 2026.}
    \item \href{}{First year into operations of LLAMAS, the Large Lenslet Array Magellan Spectrograph: design, build and commissioning overview, results and performance, lessons learned. Furesz, ... \textbf{MacBride}, ... in prep.}
    % \item  \href{}{Skipper CCD readout time optimization for astronomical applications. Wüthrich, \textbf{MacBride}, ... in prep.} 
    \item  \href{https://arxiv.org/pdf/2606.26821}{Characterizing robotic positioners under the influence of changing gravity vectors for future spectroscopic surveys. Wüthrich, ... \textbf{MacBride}, ... 25 Jun 2026.} 
    % \item \href{https://sitcomtn-148.lsst.io/}{LSST Camera Electro-Optical Test (Run 7) Results. The LSST Camera group: Antilogus, ... \textbf{MacBride}, ... 5 May 2025.}
    % \item \href{https://arxiv.org/abs/2503.07923}{The Stage-5 Spectroscopic Experiment. Besuner, Dey, ... \textbf{MacBride} ... 12 Mar 2025.}    
\end{enumerate}
\end{rSubsection}

\begin{rSubsection}{Observational and Experimental Cosmology}{}{}{}
\begin{enumerate}
    % \item \href{}{Dark Sirens in the LSST Era: Evaluating Biases and Analysis Choices. \textbf{MacBride}, Andrade-Oliveira, ... in prep.} 
    % \item \href{}{Measurement of the Hubble constant using the DECam Local Volume Exploration Survey and the fifth Gravitational-Wave Transient Catalog. \textbf{MacBride}, McMahon, ... in prep.} 
    \item \href{https://arxiv.org/pdf/2602.04766}{Measurement of the Hubble constant using the Dark Energy Survey Year 6 Gold galaxy catalog and the fourth Gravitational-Wave Transient Catalog. McMahon, ... \textbf{MacBride}, ... 5 Feb 2026.} 
    \item \href{https://arxiv.org/pdf/2512.18369}{Standard Sirens in 2040s: Probing the Cosmic Expansion History with Gravitational Waves and Spectroscopic Galaxy Surveys. Borghi, ... \textbf{MacBride}, ... 20 Dec 2025.} 
    
    
\end{enumerate}
\end{rSubsection}

\begin{rSubsection}{Multi-Messenger Astrophysics}{}{}{}
\begin{enumerate}
    % \item \href{}{The Rubin Observatory Mock-Data Challenge. \textbf{MacBride}, Chatterjee, ... in prep.} 
    % \item \href{}{Searches for Electromagnetic Counterparts of Gravitational Wave Events during LIGO-Virgo-KAGRA Observing Run 4 using the Dark Energy Camera. McMahon, \textbf{MacBride}, ... in prep.} 
    % \item \href{https://gcn.nasa.gov/circulars/43257}{LIGO/Virgo/KAGRA S251112cm: Candidates from the NSF-DOE Vera C. Rubin Observatory. Anand, \textbf{MacBride}, ...  Dec 30 2025.}
    % \item \href{https://gcn.nasa.gov/circulars/42707}{LIGO/Virgo/KAGRA S251112cm: Observations with the NSF-DOE Vera C. Rubin Observatory. \textbf{MacBride}, Anand, ...  Nov 16 2025.}
    % \item \href{https://gcn.nasa.gov/circulars/41595}{LIGO/Virgo/KAGRA S250725j: Observations with the NSF-DOE Vera C. Rubin Observatory. \textbf{MacBride}, Anand, ... August 29 2025.}
    \item \href{https://arxiv.org/abs/2508.00291}{A Joint Search for the Electromagnetic Counterpart to the Gravitational-Wave Binary Black-Hole Merger Candidate S250328ae with the Dark Energy Camera and the Prime Focus Spectrograph. Zhang, Kokubo, \textbf{MacBride}, ... 1 Aug 2025.} 
    \item \href{https://arxiv.org/pdf/2507.13409}{NSF-DOE Vera C. Rubin Observatory Observations of Interstellar Comet 3I/ATLAS (C/2025 N1) . Chandler, Bernardelli, ... \textbf{MacBride}, ... 21 July 2025.} 
    \item \href{https://arxiv.org/abs/2411.04793}{Rubin ToO 2024: Envisioning the Vera C. Rubin Observatory LSST Target of Opportunity program. Andreoni, Margutti, ... \textbf{MacBride} ... 7 Nov 2024.} 
\end{enumerate}
\end{rSubsection}

\begin{rSubsection}{Astronomical Systems Development}{}{}{}
\begin{enumerate}
    \item \href{https://arxiv.org/abs/2607.00217}{The Rubin Observatory Target-of-Opportunity Program in the First Year of Operations. \textbf{MacBride}, Ribeiro, ... 30 Jun 2026.}
    \item \ac{LSSTCam} installation and commissioning. Fanning, ... \textbf{MacBride}, ... in prep.
    \item Scheduling the LSST: commissioning the Rubin Feature Based Scheduler. Jones, ... \textbf{MacBride}, ... in prep.
    % \item \href{}{The Implementation of a Robust Observing Strategy for Discovering Optical Counterparts to Gravitational Wave Events. Fuller, \textbf{MacBride}, ... in prep.}
\end{enumerate}
\end{rSubsection}
\end{rSection}

\begin{rSection}{Outreach}
\begin{table}[h]
    \setlength{\extrarowheight}{3pt} 
\begin{tabular}{p{0.1\linewidth}p{0.8\linewidth}}
2023 - 2024 & Led group of high school students in Burlington VT to observe the 2023 annular eclipse and 2024 total eclipse. Prepared and delivered lectures for students about the physics of the sun, solar observations, and eclipses through virtual and in person presentations.\\
2023 & Designed art installation in University of Michigan Physics help room, highlighting struggles overcame by recent graduates of the physics program and success in their later lives.\\
2016 - 2020 & Showcased features of different telescopes to local tour groups at Wheaton College Observatory. Operate and adjust telescopes to show appropriate stars, objects, and events based on sky conditions. Educated children and adults in how to locate objects and explained features about the objects on sky.
    \end{tabular}
    \end{table}
% \item Physics DEI committee member from 2022 - 2024.
\end{rSection}

\begin{rSection}{Press}
\begin{table}[h]
    \setlength{\extrarowheight}{3pt} 
\begin{tabular}{p{0.1\linewidth}p{0.8\linewidth}}
    08/2025 & \href{https://www.sciencenews.org/article/vera-rubin-observatory-digital-camera}{The Vera Rubin Observatory is ready to revolutionize astronomy} - Science News. \\
    06/2025 & \href{https://www.youtube.com/watch?v=6RBZqhIgllI}{Mapping the Universe at the Rubin Observatory in Chile} - BBC. \\
    06/2025 & \href{https://www.wsj.com/science/space-astronomy/worlds-largest-digital-camera-snaps-its-first-photos-of-the-universe-68099904}{World’s Largest Digital Camera Snaps Its First Photos of the Universe} - Wall Street Journal. \\
    04/2024 & \href{https://www.vermontpublic.org/local-news/2024-04-10/nasa-volunteer-photo-eclipse-burlington-sun-atmosphere}{NASA volunteer documents eclipse from Burlington to help study sun's atmosphere} - Vermont Public Radio. \\
    04/2024 & \href{https://www.mynbc5.com/article/total-eclipse-nasa/60433342}{Vermont native and NASA-sponsored researcher returns to Burlington to study solar eclipse} \\
    % 04/2024 & \href{https://www.vermontpublic.org/local-news/2024-04-10/nasa-volunteer-photo-eclipse-burlington-sun-atmosphere}{NASA volunteer documents eclipse from Burlington to help study sun's atmosphere - Vermont Public Radio} \\
     & 
\end{tabular}   
\end{table}
\end{rSection}
\newpage
\begin{rSection}{Invited Talks}
\begin{table}[h]
    \setlength{\extrarowheight}{3pt} 
\begin{tabular}{p{0.1\linewidth}p{0.8\linewidth}}
07/2026 & Talk, Rubin Community Workshop 2026, SLAC National Accelerator Laboratory\\
07/2026 & Plenary, Rubin Community Workshop 2026, SLAC National Accelerator Laboratory\\
07/2026 & Seminar, Dark Energy School, \ac{lsst} \ac{desc} Collaboration Meeting, Paris, France\\
07/2026 & Talk, SPIE Astronomical Telescopes + Instrumentation, Copenhagen, Denmark\\
05/2026 & Talk, Rubin ToO 2026 Workshop\\
11/2025 & Talk, Wheaton College, Norton MA\\
11/2025 & Talk, Rubin-LSST France Meeting at LPSC Grenoble, France\\
09/2025 & Talk, 5th Philip Wetton Workshop, Oxford University\\
07/2025 & Talk, Rubin Community Workshop 2025, Tucson, AZ\\
% 06/2025 & Colloquium, LIGO Seminar Series\\
06/2025 & Seminar, Massachusetts Institute of Technology\\
% 05/2025 & \href{https://github.com/scimma/openMMA/wiki/Telecon20250522}{Seminar, OpenMMA Forum}\\
02/2025 & Plenary, LSST DESC Collaboration meeting\\
11/2021 & Colloquium, Wheaton College, Norton MA
\end{tabular}
\end{table}
    
\end{rSection}
\begin{rSection}{Contributed Talks}
\begin{table}[h]
    \setlength{\extrarowheight}{3pt} 
\begin{tabular}{p{0.1\linewidth}p{0.8\linewidth}}
08/2026 & Talk, Rubin 2036, SLAC National Accelerator Laboratory\\ % The Rubin Target of Opportunity Program in the First Year of LSST Operations
07/2026 & Talk, Rubin Community Workshop 2026, SLAC National Accelerator Laboratory\\ % The Rubin Target of Opportunity Program in the First Year of LSST Operations
07/2026 & Talk, Rubin Community Workshop 2026, SLAC National Accelerator Laboratory\\
07/2026 & Talk, SPIE Astronomical Telescopes + Instrumentation, Copenhagen, Denmark\\ % The Rubin Target of Opportunity Program in the First Year of LSST Operations
04/2026 & Seminar, Stanford University \\
03/2026 & Talk, LIGO-Virgo-KAGRA Collaboration Meeting, Pisa, Italy\\ % Characterization of the DECam Local Volume Exploration Survey for Standard Siren Measurement
03/2026 & Talk, LIGO-Virgo-KAGRA Collaboration Meeting, Pisa, Italy\\ % Status of the Rubin Target of Opportunity Program
05/2025 & Talk, The Dark Energy Survey Collaboration Meeting 2025, \\ % S250328ae: Joint observations with DECam and the Prime Focus Instrument at Subaru Observatory
03/2025 & Talk, APS March meeting 2025, Anaheim CA, USA\\ % LSSTCam Performance and Readiness, 
% 02/2025 & \href{https://docs.google.com/presentation/d/12EPsuWSixo40_oqrDPNaXstYnqDzhI_2aFL2hd7IwTE/edit#slide=id.g2ade3bc380_0_0}{Talk, DESC Collaboration meeting}\\
05/2020 & Talk, Wheaton College, MA\\ % Undergraduate honors thesis research
09/2018 & Seminar, Wheaton College, Norton MA\\ % REU research on dwarf satellite galaxies
08/2018 & Talk, Rutgers University\\ % REU research on dwarf satellite galaxies
\end{tabular}
\end{table}  
\end{rSection}

\end{document}
